% ==========================================================================
% GENERATED by scripts/split_paper.py. Do not edit this file.
%
% Source: imported_papers/Koopman-JEPA-TMLR-rewrite/main.tex
% Part 11: The Cayley Parameterization and Its DMD Warm-Start
%
% Edits belong upstream, in the published paper. To change what the
% chapter includes, change imports/bounded.toml and re-run
% the splitter.
% ==========================================================================
\section{The Cayley Parameterization and Its DMD Warm-Start}
\label{app:dmd}
% [orig 4.3 + orig app] DONE, verbatim, merged. The construction that used to
% sit in the methods section joins the derivations that used to sit in the
% appendix, so eq:cayley and its verification are in one place.
This appendix defines the constraint that the Koopman~(Stable) variant of Section~\ref{sec:variants} applies, and shows that it holds at every gradient step.

The soft penalty discourages $\K$ from amplifying inputs, but does not strictly prevent it. The Stable variant replaces the penalty with an architectural constraint that makes amplification impossible by construction. It forces $\K$ to be orthogonal ($\K^\top \K = I$), meaning it preserves vector norms exactly ($\|\K\z\| = \|\z\|$ for every $\z$). An orthogonal $\K$ acts as a strict rotation of the latent space, so a latent trajectory rolled out by repeated application of $\K$ can neither grow nor shrink, and divergence is structurally ruled out.

To make this constraint differentiable, we use the \emph{Cayley transform}~\citep{Erichson2019}. For an unconstrained learnable weight matrix $W \in \mathbb{R}^{D\times D}$, define the skew-symmetric matrix $A = W - W^\top$ (a matrix satisfying $A^\top = -A$), and set
\begin{equation}
  \K = (I + A)^{-1}(I - A).
  \label{eq:cayley}
\end{equation}
Numerically, we evaluate $\K$ by solving the linear system $(I+A)\K = (I-A)$ rather than computing the explicit inverse. For any $W$ the factors commute, and
\begin{equation}
    \K^\top\K
    = (I - A)^\top (I + A)^{-\top} (I + A)^{-1} (I - A)
    = (I + A) (I - A)^{-1} (I - A) (I + A)^{-1}
    = I,
\end{equation}
so $\K$ is orthogonal, and $|\lambda_i(\K)| = 1$ for every $W$ and therefore at every gradient step.

The constraint applies to the learned latent space, not to the physical state. The encoder applies a nonlinear lift rather than a change of basis, so the latent dynamics it induces need not share the spectrum of the physical system, and orthogonality of $\K$ does not assume the underlying system is lossless.

\paragraph{Initialization.}
The unconstrained and soft variants initialize $\K$ as the identity matrix perturbed by small Gaussian noise ($\sigma = 0.01$), which biases the starting eigenvalues toward the unit circle and keeps rollouts stable during early training. Since they optimize over all of $\mathbb{R}^{D\times D}$, no further care is needed. The Cayley variant is warm-started with Dynamic Mode Decomposition (DMD)~\citep{Schmid2010} instead, because training from a random $W$ was unstable in the early epochs, before the encoder has settled. DMD finds the single linear operator best matching a set of observed transition pairs, and so gives the optimizer a globally consistent starting point. The $n{=}3{,}000$ encoded transition pairs are sampled uniformly at random from the Phase 1 training buffer, and the DMD matrix is $\K_\text{DMD} = Z' Z^\dagger \in \mathbb{R}^{D\times D}$, where $Z^\dagger$ is the Moore-Penrose pseudoinverse.

\paragraph{Warm-starting the Cayley map.}
$\K_\text{DMD}$ is generally not orthogonal, so it cannot be used directly in a parameterization that requires a skew-symmetric input. We bridge the gap by inverting the Cayley map approximately and extracting its skew-symmetric component:
\begin{equation}
  A_0 = \mathrm{skew}\!\left((I + \K_\text{DMD})^{-1}(I - \K_\text{DMD})\right),
  \quad
  \mathrm{skew}(M) = \tfrac{1}{2}(M - M^\top).
\end{equation}
This initializes $W$ so that the resulting Cayley map produces an orthogonal $\K$ as close as possible to the data-driven DMD solution, which gradient optimization then refines.

% =============================================================================
